%% BirdCLEF+ 2026 working note, ceurart (PREPRINT V2 DRAFT)
\documentclass{ceurart}

\usepackage{booktabs}
\usepackage{amsmath}
\usepackage{graphicx}
\usepackage{multirow}
\usepackage{hyperref}

\begin{document}

\copyrightyear{2026}
\copyrightclause{Copyright for this paper by its authors.
  Use permitted under Creative Commons License Attribution 4.0 International (CC BY 4.0).}

\conference{CLEF 2026: Conference and Labs of the Evaluation Forum,
  September 21--24, 2026, Jena, Germany}

\title{Partial-Label Evaluation of Prototype Networks with Temporal Aggregation in BirdCLEF+ 2026}

\author[1]{Tsz Chai Alan Liu}[%
  email=tcaliu@sonoform.org,
]
\address[1]{Sonoform Labs, Ontario, Canada}

%% ============================================================
\begin{abstract}
BirdCLEF+ 2026 requires identifying 234 bird, amphibian, insect, mammal, and reptile species from Brazilian Pantanal soundscape recordings under a strict 90-minute CPU-only inference deadline. Prototype-based models are attractive in this setting because they can attach predictions to interpretable acoustic evidence, but their offline evaluation may cover only a subset of the deployment label space. We study this problem by using a fixed Perch+prototype+SED rank-fusion pipeline as a deployment testbed and replacing only its prototype member with author-built prototype heads trained on frozen Perch embeddings. The deployed prototype member uses temporal aggregation across soundscape windows, while the primary replacement heads make window-level predictions. On a 708-window labelled soundscape bank covering 71 of 234 species, the replacement heads form an interpretability--accuracy frontier, and the higher-accuracy D2 head outperforms the deployed prototype member on the covered-species offline endpoint. However, this ranking does not transfer to hidden-test deployment: D1 scores 0.938 private macro-AUC, D2 scores 0.937, and a raw-Perch version without the prototype member scores 0.939, while the deployed temporally aggregated prototype member remains best at 0.942. Mechanism diagnostics indicate that the replacement heads degrade ranking on the 163 species absent from the offline bank. A no-SSM ablation reduces the deployed head from 0.94218 to 0.94089 private macro-AUC, showing that temporal aggregation contributes to the advantage but does not fully explain it. A supporting matched-control study shows that prototype explanations can be class-pure and call-focused in a cleaner setting, so the result is not simply a failure of prototypes as explanations. Overall, our experiments document a scoped failure mode of partial-label evaluation: a prototype component can improve a covered-species offline benchmark while reducing performance in the full bioacoustic deployment system.

\end{abstract}

\begin{keywords}
  BirdCLEF+ 2026 \sep
  Bioacoustics \sep
  Prototype-based Models \sep
  Temporal Aggregation \sep
  Partial-Label Evaluation \sep
  Offline--Deployment Gap
\end{keywords}

\maketitle

%% ============================================================
%%  1. INTRODUCTION
%% ============================================================
\section{Introduction}

The BirdCLEF+ 2026 challenge~\cite{birdclef2026}, part of LifeCLEF 2026~\cite{lifeclef2026}, requires identifying 234 bird, amphibian, insect, mammal, and reptile species from Brazilian Pantanal soundscape recordings under a 90-minute CPU-only inference budget. The task extends the BirdCLEF series from bird-only classification toward broader passive-acoustic biodiversity monitoring~\cite{kahl2024birdclef,kahl2025birdclef}. This setting introduces several practical challenges:

\begin{itemize}
\item Significant domain shift between focal training recordings and passive-acoustic-monitoring soundscapes~\cite{ganin2015domain};
\item Severe class imbalance, with some common species having hundreds of recordings and rare species having fewer than ten;
\item Partial label coverage in offline soundscape evaluation, where the labelled soundscape bank used in this work contains positives for only 71 of 234 target species;
\item Incomplete foundation-model coverage, with 28 of 234 species lacking Perch~\cite{ghani2023global} logit mappings.
\end{itemize}

Many strong BirdCLEF systems address these constraints by combining pretrained audio models, prototype or embedding-based components, sound-event-detection models, rank fusion, and post-processing. In this work, we use a fixed ProtoSSMv2+SED rank-fusion core, implemented from public Kaggle reference notebooks~\cite{protossmv2}, as a deployment testbed. This fixed core scored 0.942 private macro-AUC. We hold the core fixed and replace only its prototype member with author-built prototype heads trained on frozen Perch embeddings.

This controlled replacement setup lets us study a narrow evaluation question: when prototype components are evaluated offline on a partially covered soundscape bank, do their offline rankings survive deployment inside the full 234-species system? The question is especially relevant because the deployed prototype member uses temporal aggregation across soundscape windows, whereas our primary replacement heads are simpler window-level prototype models. Thus, the comparison tests not only prototype-head design, but also whether offline window-level evaluation is sufficient for judging a temporally aggregated component inside a full competition system.

Our results show that offline evaluation can be misleading in this setting. On the 708-window labelled soundscape bank covering 71 species, the author-built heads form an interpretability--accuracy frontier, and the higher-accuracy D2 head exceeds the deployed prototype member on the covered-species offline endpoint. However, when the heads are swapped into the full system and scored on the hidden test, neither primary replacement improves over dropping the prototype member for raw Perch logits. D1 scores 0.938 private, D2 scores 0.937, raw Perch scores 0.939, and the deployed temporally aggregated prototype member remains best at 0.942. Diagnostics indicate that the failure is concentrated on the 163 species absent from the offline bank: the replacement heads degrade ranking on species the bank never evaluates, while the deployed prototype member preserves stronger signal for them. A no-SSM ablation reduces the deployed head to 0.94089 private macro-AUC, indicating that temporal aggregation contributes to the deployed advantage but does not fully explain it.

Our contributions are:

\begin{itemize}
\item A controlled system-level replacement protocol that swaps only the prototype member of a fixed Perch+prototype+SED rank-fusion pipeline;
\item An offline interpretability--accuracy frontier for author-built prototype heads trained on frozen Perch embeddings;
\item Hidden-test evidence that the offline ranking does not transfer to the full 234-species system;
\item A mechanism analysis showing that the failure is concentrated on species absent from the offline soundscape bank;
\item A temporal-module ablation showing that temporal aggregation contributes to, but does not fully explain, the deployed prototype member's advantage.
\end{itemize}

A supporting matched-control study shows that prototype explanations can be class-pure and call-focused in a cleaner setting, suggesting that the observed failure is not simply a failure of prototypes as explanations. Instead, we frame it as a failure mode of partial-label evaluation for interpretable components in BirdCLEF-style systems.

The remainder of this paper is organized as follows. Section~\ref{sec:related} reviews related work. Section~\ref{sec:data} describes the task, dataset, and label-coverage gap. Section~\ref{sec:system} presents the fixed deployment testbed. Section~\ref{sec:part2} describes the replacement prototype heads and offline results. Section~\ref{sec:leaderboard} reports hidden-test replacement results. Section~\ref{sec:mechanism} analyzes the uncovered-species failure mode. Section~\ref{sec:validity} summarizes the supporting matched-control study. Sections~\ref{sec:discussion}--\ref{sec:future} discuss implications, limitations, and future work.

%% ============================================================
%%  2. RELATED WORK
%% ============================================================
\section{Related Work}
\label{sec:related}

\subsection{Prototype networks and interpretability}

Prototype networks aim to make classification decisions interpretable by comparing an input to learned representative examples. ProtoPNet~\cite{chen2019protopnet} introduced this paradigm for image classification: a convolutional backbone extracts latent features, a prototype layer measures similarity to learned class-specific prototypes, and the final layer maps prototype activations to class predictions. This produces a ``this looks like that'' form of explanation in which the prototypes are part of the model rather than a post-hoc visualization.

However, prototype interpretability is not automatic. Prior work has shown that latent prototype similarity may fail to correspond to meaningful input-space similarity, especially when prototypes capture artifacts or compressed features rather than semantically relevant structure~\cite{hoffmann2021looks}. This motivates evaluating prototype quality quantitatively rather than relying only on visual or acoustic examples. In this work, we measure prototype validity using class purity and shortcut rate in a matched-control setting, and then test whether prototype components that look strong offline remain useful inside a deployed BirdCLEF system.

\subsection{Prototype and foundation models in bioacoustics}

AudioProtoPNet~\cite{heinrich2025audioprotopnet} adapts the prototype-network paradigm to multi-label bird sound classification, using a ConvNeXt backbone pretrained on BirdSet~\cite{rauch2024birdset}. It shows that prototype explanations can be useful in bird-only bioacoustic classification, but leaves broader multi-taxon settings as future work. Our supporting matched-control study probes this boundary by testing whether prototype explanations remain class-pure and call-focused in a cleaner audio setting before studying their failure inside the full BirdCLEF+ 2026 deployment system.

Pretrained bioacoustic models have become central to recent BirdCLEF systems. BirdNET~\cite{kahl2021birdnet}, Perch~\cite{ghani2023global}, PANNs~\cite{kong2020panns}, and BirdSet backbones~\cite{rauch2024birdset} provide strong feature extractors for downstream classification. In BirdCLEF+ 2026, we use frozen Perch embeddings and logits as the fixed representation layer for all prototype-head replacement experiments. This lets us study the prototype component itself while holding the backbone representation constant.

\subsection{BirdCLEF systems and component evaluation}

Recent BirdCLEF working notes emphasize transfer learning, pseudo-labeling, distillation, ensembling, and post-processing as practical tools for bioacoustic classification under domain shift and strict inference budgets. Miyaguchi et al.~\cite{miyaguchi2025stsg} explored lightweight tokenized spectrogram embeddings for fast inference, while Sydorskyi and Gon\c{c}alves~\cite{sydorskyi2025domainshift} used semi-supervised distillation and iterative pseudo-labeling to achieve a private AUC of 0.928 in BirdCLEF+ 2025. These studies report system-level performance and ablations, but they do not focus on whether an offline ranking of an interpretable component transfers when that component is inserted into the full deployment system.

This component-versus-system question is important for BirdCLEF+ 2026 because offline soundscape evaluation may cover only part of the deployment label space. In our case, the labelled soundscape bank used for offline prototype evaluation contains positives for only 71 of 234 target species. A prototype head can therefore improve the covered-species offline benchmark while damaging ranking on species the offline bank cannot evaluate. Our work studies this failure mode directly by replacing only the prototype member of a fixed Perch+prototype+SED rank-fusion pipeline, measuring the effect on the hidden 234-species leaderboard, and adding a temporal-module ablation to separate partial-label effects from the deployed member's temporal aggregation.


%% ============================================================
%%  3. TASK AND DATA
%% ============================================================
\section{Task and Data}
\label{sec:data}

\subsection{Task overview}

BirdCLEF+ 2026~\cite{birdclef2026} targets 234 species from the Brazilian Pantanal, expanding beyond previous editions' bird-only scope~\cite{kahl2024birdclef, kahl2025birdclef}.
Participants predict species-presence probabilities for each 5-second window of one-minute soundscape recordings.
The evaluation metric is macro-averaged ROC-AUC, skipping classes with no true positives in the test set\footnote{\url{https://www.kaggle.com/code/metric/birdclef-roc-auc}}.
All inference must complete within a 90-minute CPU-only budget on a Kaggle notebook (Intel Xeon, 4 cores).

\subsection{Training data and species distribution}

\begin{table}[t]
\caption{Species distribution, approximate focal recording counts, and direct Perch logit coverage by taxon. Proxy mappings are described separately in Section~\ref{sec:data}.}
\label{tab:taxa}
\begin{tabular}{lrrr}
\toprule
Taxon & Species & Focal rec. & Direct Perch \\
\midrule
Aves       & 162 & $\sim$24,000 & 159 \\
Amphibia   & 35  & $\sim$1,800  & 22 \\
Insecta    & 28  & $\sim$900    & 13 \\
Mammalia   & 8   & $\sim$300    & 6 \\
Reptilia   & 1   & $\sim$20     & 0 \\
\midrule
\textbf{Total} & \textbf{234} & $\sim$\textbf{27,000} & \textbf{203} \\
\bottomrule
\end{tabular}
\end{table}

Table~\ref{tab:taxa} summarizes the dataset composition.
Training data includes focal recordings from three sources: Xeno-Canto\footnote{\url{https://xeno-canto.org/}}, iNaturalist\footnote{\url{https://www.inaturalist.org/}}, and the Colombian Sound Archive (CSA).
Labelled soundscapes with 5-second multi-label annotations are also provided.
Class imbalance is severe: the most-represented bird species have hundreds of recordings, while several mammal and insect species have fewer than 10.
This imbalance is central to our study: the offline evaluation bank covers only 71 of 234 species with positive instances, and that coverage gap drives our main finding (Section~\ref{sec:mechanism}).

\subsection{Domain shift}

The core difficulty is a focal$\rightarrow$soundscape domain shift, well-documented in prior BirdCLEF editions~\cite{sydorskyi2025domainshift}.
Focal recordings are clean, single-species clips recorded at close range with directional microphones.
Soundscapes are captured by passive acoustic monitoring devices, contain overlapping species (approximately 4.2 per positive window), background noise from rain and wind, and varying recording conditions (Figure~\ref{fig:domainshift}).

\begin{figure}[t]
\centering
\includegraphics[width=\columnwidth]{fig_domain_shift.pdf}
\caption{Focal$\rightarrow$soundscape domain shift for one species (Dwarf Tree Frog): a clean, single-species focal recording (top) vs.\ a soundscape window with eight overlapping species and environmental noise (bottom). This pair is illustrative rather than a quantitative summary of all recordings. The overlapping vocalizations and broadband noise absent from focal clips are the primary challenge for both classification and prototype projection.}
\label{fig:domainshift}
\end{figure}

This shift provides deployment context for our experiments. However, the main failure analyzed below is not domain shift alone, but the interaction between partial offline label coverage and full-system deployment.

\subsection{Perch coverage gap}

A separate coverage issue comes from the Perch class mapping. Of the 234 target species, 203 have direct Perch logit mappings, 3 are handled through taxonomic proxy mappings, and 28 produce no Perch logit signal. Table~\ref{tab:taxa} reports only direct Perch coverage by taxon to avoid conflating target species with proxy logits. The unmapped species are concentrated mainly among insects and amphibians.

This Perch coverage gap is distinct from the offline evaluation gap that drives our main mechanism. The Perch gap concerns which species receive foundation-model logit support, whereas the offline evaluation gap concerns which species have positive labels in the 708-window soundscape bank. In our offline bank, only 71 of 234 species have positive examples; the remaining 163 species cannot contribute to offline macro-AUC. Section~\ref{sec:mechanism} shows that the deployment advantage of the original prototype member is concentrated on these offline-absent species.

\subsection{Offline soundscape bank and split construction}
\label{sec:evalbank}

For local prototype development and offline evaluation, we use a labelled soundscape bank drawn from the competition-provided 5-second multi-label soundscape annotations for which Perch embeddings and logits were available in our pipeline. The bank contains 708 windows from 59 one-minute source files. Since each one-minute file is divided into twelve 5-second windows, each source file contributes its temporal windows as a group. We do not rebalance the bank by species or taxon.

All offline folds are constructed at the source-file level using GroupKFold(5). In each fold, all windows from a given one-minute source file are assigned entirely to either the training split or the held-out split; no windows from the same source file cross the train--validation boundary. This prevents temporally adjacent windows from the same recording from appearing on both sides of an offline comparison. Species are included in offline macro-AUC only if they have at least one positive instance in the 708-window bank, yielding 71 evaluable species. The remaining 163 target species have no positive examples in this bank, so the offline endpoint can evaluate covered-species behavior but cannot evaluate whether a head damages ranking for absent species. This coverage limitation becomes central in Section~\ref{sec:mechanism}.

The primary D1 and D2 replacement heads and the offline design sweeps in Tables~\ref{tab:buildup}--\ref{tab:pertaxon} are developed on this soundscape bank. Focal recordings are not used to train the primary D1 and D2 heads; they are introduced only in the separate D2-FOCAL sensitivity experiment (Section~\ref{sec:mechanism}). Thus, the main offline endpoint and the primary replacement heads share the same labelled soundscape source, while the hidden leaderboard provides a separate system-level deployment check.

\subsection{Evaluation endpoints and analysis plan}
\label{sec:endpoints}

We distinguish component-level offline evaluation from system-level deployment evaluation. The offline endpoint is out-of-fold macro-AUC on the 708-window labelled soundscape bank, computed over the 71 species with at least one positive instance. This endpoint is used to compare prototype-head designs on the covered-species bank and to construct the interpretability--accuracy frontier. The deployment endpoint is private hidden-test macro-AUC after replacing the prototype member inside the fixed ProtoSSMv2+SED core. The primary deployment comparison is between the deployed ProtoSSMv2 member, the two primary replacement heads (D1 and D2), and the raw-Perch drop baseline.

We treat the other analyses as exploratory diagnostics or sensitivity checks. Coverage-gating, re-weighting, D2-FOCAL, per-taxon breakdowns, absent-species diagnostics, and the supporting AudioProtoPNet matched-control study are not used as confirmatory endpoints; they are used to interpret the direction of the deployment reversal and to probe possible mechanisms. Because the local bank is small and the hidden-test labels are unavailable, we present the paper as a documented system-level case study rather than as a fully powered hypothesis test.

\subsection{Statistical reporting conventions}
\label{sec:statreporting}

For offline experiments, reported macro-AUC values are out-of-fold estimates from the file-level GroupKFold procedure. Fold-level summaries, when available, are used as descriptive uncertainty summaries rather than as independent replicated experiments, because the number of folds is small and species coverage is incomplete. Hyperparameter sweeps on the offline bank are treated as model-development analyses; they should not be read as a locked, independent model-selection benchmark.

Prototype presence-purity is reported over the set of evaluated prototypes. Null purity values are computed by shuffling prototype-to-class assignments. When confidence intervals are reported for purity, they summarize uncertainty over prototypes rather than hidden-test deployment uncertainty. For hidden-test leaderboard results, Kaggle returns only aggregate public/private macro-AUC scores and does not expose per-example or per-species labels. We therefore do not compute paired tests, per-species confidence intervals, or p-values for leaderboard differences. Public and private scores are reported as descriptive system-level outcomes, and no multiple-comparison correction is applied to exploratory subgroup or rescue analyses.


%% ============================================================
%%  4. COMPETITION SYSTEM
%% ============================================================
\section{Deployment Testbed}
\label{sec:system}

We use a fixed ProtoSSMv2+SED rank-fusion core as the deployment testbed for all replacement experiments. This core contains two active members: the deployed ProtoSSMv2 prototype member and a Perch-distilled SED member. Auxiliary EfficientNet members from the original public pipeline are disabled uniformly across the replacement experiments, so only the prototype member differs between leaderboard rows. Each leaderboard row in Table~\ref{tab:leaderboard} corresponds to one submitted system generated by replacing the prototype-head artifact and leaving the rest of this core unchanged.

Each member's scores are converted to per-species percentile ranks and then weighted-summed. For member $m$, species $c$, raw score $s_c^{(m)}$, rank-percentile score $r_c^{(m)}$, and fusion weight $w_m$, the final fused score is
\[
  r_c^{(m)} = \operatorname{rankpercentile}(s_c^{(m)}), \qquad
  S_c = \sum_m w_m r_c^{(m)}.
\]
Rank-space blending is natural for macro-AUC because the metric depends only on within-species ordering. The deployed ProtoSSMv2+SED core scores 0.942 private macro-AUC on the hidden test. Since hidden labels are not released, this score is used only as the deployment reference for descriptive replacement comparisons.

The prototype member, ProtoSSMv2~\cite{protossmv2}, uses one class-specific prototype per species, projected to a 320-dimensional space. It computes cosine similarity between Perch-derived features and class prototypes, applies positive-temperature scaling via softplus, and blends prototype scores with Perch logits using a learned per-class fusion gate $\gamma_c = \sigma(\alpha_c)$. A bidirectional residual selective state-space model~\cite{gu2024mamba,gu2022s4} aggregates predictions across the 12 temporal windows in each one-minute soundscape file.


The SED model uses an EfficientNet~\cite{tan2019efficientnet} backbone distilled from Perch embeddings (MSE loss), producing clip-level and frame-level logits combined as $p = 0.5 \cdot \sigma(l_\text{clip}) + 0.5 \cdot \sigma(\max_t l_\text{frame}(t))$, averaged across 5 folds with Gaussian smoothing ($\sigma{=}0.65$).

Post-processing includes empirical-Bayes ecological priors (site/hour shrinkage estimators with circular-Gaussian hour smoothing, $\sigma{=}1.5$ hours) and rank-active methods (taxonomy smoothing, file-level max pooling, top-2 averaging, and $\Delta$-smoothing).
We verified that per-class temperature scaling and isotonic thresholding are rank-inert (0.00 rank change), confirming they cannot affect AUC.

We do not claim novelty for the base pipeline; it is used as a fixed deployment testbed. Our contribution is the controlled replacement of its prototype member and the analysis around that replacement. Table~\ref{tab:ownership} states what is public and what is ours.

\begin{table}[t]
\caption{What is public versus ours. The deployed pipeline is a fixed, public testbed; our contribution is the controlled replacement of its prototype member and the analysis around it. D1 denotes the pure multi-prototype replacement head; D2 denotes the input-gated high-offline-AUC replacement head; eco. priors denotes empirical-Bayes ecological priors.}
\label{tab:ownership}
\begin{tabular}{lll}
\toprule
Component & Source & Role in this study \\
\midrule
Perch v2 embeddings / logits & public model~\cite{ghani2023global} & fixed testbed \\
ProtoSSMv2 prototype member & public notebook~\cite{protossmv2} & fixed testbed (replaced) \\
Perch-distilled SED member & public notebook~\cite{protossmv2} & fixed testbed \\
Rank fusion + eco.\ priors & public notebook~\cite{protossmv2} & fixed testbed \\
D1 / D2 / static-gate heads & \textbf{this work} & \textbf{intervention} \\
Controlled replacement + analysis & \textbf{this work} & \textbf{contribution} \\
Matched-control study (App.~\ref{app:partone}) & \textbf{this work} & \textbf{contribution} \\
\bottomrule
\end{tabular}
\end{table}


%% ============================================================
%%  6. PART II: AUTHOR-BUILT HEADS
%% ============================================================
\section{Prototype Head Design and Offline Evaluation}
\label{sec:part2}

\subsection{Replacement-head setup}
\label{sec:headsetup}

We design and train replacement prototype classification heads on frozen Perch v2 embeddings (1,536-d). For an input window $x$, let $z(x)$ denote the frozen Perch embedding and $\ell_c^{\mathrm{Perch}}(x)$ denote the Perch logit for species $c$ when available. Each prototype head contains one or more learned prototypes $p_{c,k}$ assigned to species $c$. With normalized embeddings and prototypes, the prototype activation is
\[
  a_{c,k}(x)=\left\langle \frac{z(x)}{\|z(x)\|},\frac{p_{c,k}}{\|p_{c,k}\|}\right\rangle,
  \qquad
  a_c(x)=\max_k a_{c,k}(x).
\]
The prototype logit is then
\[
  h_c(x)=\operatorname{softplus}(\eta_c)\,a_c(x)+b_c,
\]
where $\eta_c$ is the learnable temperature parameter and $b_c$ is a per-class bias. For pure prototype heads, $h_c(x)$ is the class score. For fusion heads, the prototype logit is combined with the Perch logit through a convex gate,
\[
  q_c(x)=\gamma_c(x)h_c(x)+(1-\gamma_c(x))\ell_c^{\mathrm{Perch}}(x).
\]
The static-gate head uses $\gamma_c=\sigma(\alpha_c)$, while D2 uses an input-dependent gate $\gamma_c(x)$. D1 denotes the pure multi-prototype replacement head with interpretability losses; D2 denotes the input-gated high-offline-AUC replacement head.

Training uses FocalBCE loss~\cite{lin2017focal}, AdamW~\cite{loshchilov2019adamw}, and cosine annealing over 60 epochs. When interpretability losses are active, the training objective is
\[
  \mathcal{L}=\mathcal{L}_{\mathrm{FocalBCE}}+\lambda_{\mathrm{clust}}\mathcal{L}_{\mathrm{clust}}+\lambda_{\mathrm{sep}}\mathcal{L}_{\mathrm{sep}}.
\]
Using cosine distance $d(u,p)=1-\langle u/\|u\|,p/\|p\|\rangle$, the cluster loss pulls each prototype toward windows containing its assigned class, while the separation loss penalizes closeness to confirmed-absent windows for that class. The multi-label-aware separation variant changes only the negative set used by $\mathcal{L}_{\mathrm{sep}}$: in multi-species windows, co-occurring positive species are not treated as negatives for one another.

\subsection{Training folds and model selection}
\label{sec:modelselection}

For the primary replacement heads and the offline sweeps, each GroupKFold split is trained independently on four file groups and evaluated on the held-out file group. Out-of-fold predictions are concatenated across folds before computing macro-AUC over the 71 covered species. The offline design sweeps are model-development analyses on the same local bank: D2 was selected as the high-offline-AUC replacement head, while D1 was selected as the high-purity pure-prototype comparison. Because the same local bank guided design choices, the offline frontier should be interpreted as a descriptive development result rather than as a completely independent final test. The hidden-test replacement submissions provide the separate deployment check, but those scores are also descriptive because hidden labels are unavailable.

Perch embeddings are strongly species-discriminative: a t-SNE projection of focal recordings (Figure~\ref{fig:embeddings}) shows that different species occupy well-separated regions of the embedding space, which is what makes a single prototype per species a viable starting point.

\begin{figure}[t]
\centering
\includegraphics[width=\linewidth]{fig_embeddings.pdf}
\caption{t-SNE projection of frozen Perch~v2 embeddings for focal recordings of twelve representative species (eBird/iNat codes); each point is one recording window. Axes are arbitrary t-SNE coordinates. Distinct per-species clusters show that the frozen embedding space is highly species-discriminative, which is the basis for single-prototype-per-species heads. Evaluation instead uses noisier multi-species soundscape windows (Section~\ref{sec:mechanism}).}
\label{fig:embeddings}
\end{figure}
Prototype \textbf{presence-purity} is defined as: for each prototype, whether the nearest soundscape window in the 708-window bank (by cosine similarity) contains the prototype's assigned class. For prototype set $\mathcal{P}$, assigned class $c(p)$, and nearest window $x_p^\star$, we compute
\[
  \operatorname{purity}
  = \frac{1}{|\mathcal{P}|}\sum_{p\in\mathcal{P}}
  \mathbf{1}\{y_{c(p)}(x_p^\star)=1\}.
\]

\subsection{Deployed head audit}

Before building new heads, we audited the deployed ProtoSSMv2~\cite{protossmv2} to establish the reference we aim to match.

\paragraph{Architecture.}
ProtoSSMv2 uses 234 class-specific prototypes (one per species, 320-d after internal projection), a positive-temperature diagonal cosine readout, a per-class fusion gate $\gamma_c = \sigma(\alpha_c)$, and a bidirectional residual SSM~\cite{gu2024mamba}.

\paragraph{Fusion gate analysis.}
The learned gate values are flat: $\gamma \approx 0.50$ across all 234 species and all taxa (global std 0.003, range [0.493, 0.507]).
Table~\ref{tab:gate_taxon} confirms the uniformity.
The gate did not differentiate prototype vs.\ Perch reliance in this run, a negative result.
This stands in contrast to the design motivation (graceful coverage degradation for unmapped species) and suggests the gate may require longer training or stronger initialization to differentiate.

\begin{table}[t]
\caption{ProtoSSMv2 fusion gate $\gamma$ by taxon. The gate is uniform everywhere, including for unmapped species.}
\label{tab:gate_taxon}
\begin{tabular}{lrr}
\toprule
Taxon & Mean $\gamma$ & Std \\
\midrule
Amphibia  & 0.497 & 0.003 \\
Aves      & 0.501 & 0.003 \\
Insecta   & 0.503 & 0.003 \\
Mammalia  & 0.499 & 0.002 \\
Reptilia  & 0.501 & -- \\
\bottomrule
\end{tabular}
\end{table}

\paragraph{Prototype validity.}
ProtoSSMv2's prototype purity is 0.789 [95\% CI 0.69--0.87], 42$\sigma$ above the permutation null (0.024 $\pm$ 0.018).
Interestingly, the per-taxon pattern runs opposite to AudioProtoPNet's gradient:
Insecta 0.96 $>$ Amphibia 0.88 $>$ Aves 0.64 (Mammalia 0.33, $n{=}3$; Reptilia 0.00, $n{=}1$; both quarantined); see Figure~\ref{fig:taxon_purity}.

\begin{figure}[t]
\centering
\includegraphics[width=\columnwidth]{fig_taxon_purity.pdf}
\caption{Per-taxon prototype purity for AudioProtoPNet (matched-control study, Appendix~\ref{app:partone}; ConvNeXt backbone) vs.\ ProtoSSMv2 (Perch embeddings). The gradients are opposite: AudioProtoPNet is bird-clean; ProtoSSMv2 is insect-clean. The contrast is consistent with prototype validity tracking training-domain acoustic properties rather than the prototype mechanism per se, but the comparison is confounded by backbone, training data, and prototype granularity. Insect purity for AudioProtoPNet is coverage-limited; see text.}
\label{fig:taxon_purity}
\end{figure}

This pattern survives a frequency control.
There is a positive correlation between purity and log-sample-count (Spearman $\rho = 0.30$), but a logistic regression with the log-count covariate retains the taxon effect: Insecta $+1.40$, Aves $-0.71$ (vs.\ Amphibia reference).
Insect prototypes project cleanly even when the species is rare (purity 1.0 at low sample count), while avian prototypes are frequency-dependent (0.47 for rare birds, 1.0 for common birds).
This pattern is consistent with validity tracking acoustic distinctiveness: insects produce spectrally unique calls that are easier to separate, while birds in a dense avian community overlap spectrally and may require more data to differentiate.

\subsection{The offline frontier}

\begin{table}[t]
\caption{Stepwise build-up. Each row adds one component to a prototype head on frozen Perch embeddings. $\Delta$ values relative to the previous row. Bold = best in column.}
\label{tab:buildup}
\begin{tabular}{rlcccc}
\toprule
\# & Config & AUC & $\Delta$ & Pur. & $\Delta$ \\
\midrule
-- & Perch alone         & .739 & --     & --  & -- \\
1  & Linear probe        & .766 & +.027   & --  & -- \\
2  & 1-proto, no loss    & .762 & $-.004$ & .69  & -- \\
3  & + cluster           & .604 & $-.158$ & .944 & +.254 \\
4  & + separation        & .605 & +.001   & .958 & +.014 \\
5  & + static gate       & .742 & +.137   & .958 & .000 \\
6  & + input gate (D2)   & \textbf{.830} & +.088 & .437 & $-.521$ \\
7  & 3-proto+loss (D1)   & .606 & --     & \textbf{.958} & -- \\
-- & ProtoSSMv2          & .803 & --     & .790 & -- \\
\bottomrule
\end{tabular}
\end{table}

\begin{figure}[t]
\centering
\includegraphics[width=\columnwidth]{fig_frontier.pdf}
\caption{The interpretability--accuracy frontier. Each point is a head configuration from Table~\ref{tab:buildup}. Cluster loss increases purity at steep accuracy cost (rows 2$\rightarrow$3). The input-dependent gate recovers accuracy while reducing purity (rows 5$\rightarrow$6). Row~5 (static gate + losses) is the only design point in this sweep where interpretability is approximately free. The deployed ProtoSSMv2 sits in the interior of the frontier.}
\label{fig:frontier}
\end{figure}

Table~\ref{tab:buildup} and Figure~\ref{fig:frontier} reveal a clean interpretability--accuracy frontier.
We highlight three findings:

\paragraph{Cluster loss is the interpretability lever.}
Adding cluster loss (rows 2$\rightarrow$3) buys $+0.254$ purity at $-0.158$ AUC, the single largest effect in the table.
Separation loss adds only $+0.014$ marginal purity (rows 3$\rightarrow$4).
We also tested a multi-label-aware separation variant that penalizes proximity only to confirmed-absent species (relevant for soundscapes with co-occurring species that should not be pushed apart); it performs similarly to standard separation (AUC 0.605 vs.\ 0.603), a negative result.
Increasing prototype count from 1 to 3 (row~7) does not help: capacity is not the bottleneck.

\paragraph{The input-dependent gate trades purity for accuracy.}
Replacing the static gate with an input-dependent gate (rows 5$\rightarrow$6) gains $+0.088$ AUC at $-0.521$ purity.
The input gate increases reliance on Perch logits and reduces the prototype's discriminative role.
We verified this tension is fundamental by building a ``full-spec'' D2 with both interpretability losses \emph{and} the input gate: it scores 0.691 / 0.958.
The accuracy collapses below even the static gate.
In these sweeps, no gated configuration achieves both high accuracy and high purity.

\paragraph{Row~5 is the free-interpretability sweet spot.}
A static-gate fusion with cluster and separation losses (row~5) recovers Perch-level AUC (0.742) while maintaining 0.958 purity, the only design point where interpretability is roughly free.
However, its absolute accuracy is far below D2 (0.830) or the deployed ProtoSSMv2 (0.803), limiting its practical utility.

\subsection{Design sweeps}

\begin{table}[t]
\caption{Design sweeps. $\Delta$ values are relative to each group's baseline (shown in parentheses in the group header). Pur. denotes prototype presence-purity; a dash indicates that the quantity was not measured for that sweep or is not applicable.}
\label{tab:sweeps}
\begin{tabular}{llrr}
\toprule
Sweep & Value & $\Delta$AUC & $\Delta$purity \\
\midrule
\multicolumn{4}{l}{\emph{Pure-proto head (base: .605 / .958)}} \\
\midrule
\multirow{2}{*}{Prototype count $K$}
  & $K{=}3$         & +.001 & .000 \\
  & $K{=}10$        & +.002 & $-.518$ \\
\midrule
Projection  & linear 320-d    & +.042 & .000 \\
\midrule
Separation loss  & multi-label     & $-.002$ & .000 \\
\midrule
\multicolumn{4}{l}{\emph{Gated head (base: .830 / .437)}} \\
\midrule
\multirow{2}{*}{Prototype count $K$}
  & $K{=}3$         & +.003 & +.105 \\
  & $K{=}5$         & +.006 & +.169 \\
\midrule
Epochs      & 25 (peak)       & +.018 & -- \\
\midrule
Gate type   & static scalar   & $-.088$ & +.521 \\
\bottomrule
\end{tabular}
\end{table}

Table~\ref{tab:sweeps} reports additional design sweeps.
Two contrasting patterns emerge:
on the pure-prototype head, increasing $K$ leaves AUC flat while per-prototype purity collapses (0.958 to 0.440), as extra prototypes dilute discriminative signal;
on the gated head, $K$ mildly improves both AUC ($+0.006$ at $K{=}5$) and purity ($+0.169$), because the gate can selectively attend to the best prototype.

A linear projection from 1,536-d to 320-d provides a purity-preserving AUC gain on the pure head ($+0.042$), suggesting dimensionality reduction removes noise that hurts classification without affecting prototype alignment.
D2 training peaks at 25 epochs (AUC 0.848) and overfits beyond, consistent with the small training bank (708 windows).
Temporal smoothing across windows was neutral to negative.

\subsection{Per-taxon accuracy}

\begin{table}[t]
\caption{Per-taxon OOF macro-AUC for key configs (covered species only; Mammalia/Reptilia omitted as rows but included in All). Bold = best per row.}
\label{tab:pertaxon}
\begin{tabular}{lrrr}
\toprule
Taxon & Linear & D2 (gated) & D1 (pure) \\
\midrule
Aves (25)      & .712 & \textbf{.804} & .552 \\
Amphibia (17)  & .799 & \textbf{.884} & .687 \\
Insecta (25)   & .820 & \textbf{.827} & .619 \\
All (71)       & .766 & \textbf{.830} & .606 \\
\bottomrule
\end{tabular}
\end{table}

Table~\ref{tab:pertaxon} shows that D2's accuracy advantage over the linear probe is largest where Perch is weakest (Aves $+0.092$, Amphibia $+0.085$) and near zero on Insecta ($+0.007$), where Perch embeddings are already discriminative.
D1's accuracy collapse is taxon-general.
The accuracy-gain gradient is orthogonal to the purity gradient (Insecta-cleanest), confirming that accuracy and interpretability respond to different properties of the data.


%% ============================================================
%%  7. LEADERBOARD: THE INVERSION
%% ============================================================
\section{Leaderboard: The Offline Ranking Reverses}
\label{sec:leaderboard}

\begin{table}[t]
\caption{Leaderboard evaluation. Each row replaces or ablates the prototype member in the deployed Proto+SED core and is scored on the hidden test. Auxiliary CNN members are disabled uniformly. D1 is the pure prototype head; D2 is the input-gated head; no-SSM ablates the deployed head's temporal module; cov-gated uses prototype scores only for species covered by the offline bank; Drop uses raw Perch logits in place of the prototype member. $\Delta$ = private macro-AUC minus the ProtoSSMv2 baseline; negative values are worse.}
\label{tab:leaderboard}
\begin{tabular}{lccrc}
\toprule
Head & Off. & Pub. & Priv. & $\Delta$ \\
\midrule
ProtoSSMv2          & .803 & .949 & \textbf{.942} & -- \\
ProtoSSMv2 no-SSM   & --   & .948 & .941 & $-.001$ \\
D1 cov-gated        & .606 & .944 & .940 & $-.003$ \\
ProtoSSMv2 cov-g.   & .803 & .946 & .940 & $-.003$ \\
D2 re-weighted      & .830 & .942 & .939 & $-.003$ \\
Drop (Perch)        & .739 & .945 & .939 & $-.003$ \\
D1 pure             & .606 & .938 & .938 & $-.004$ \\
D2 cov-gated        & .830 & .944 & .938 & $-.005$ \\
D2 input-gated      & \textbf{.830} & .941 & .937 & $-.005$ \\
D2-FOCAL            & --  & .939 & .935 & $-.007$ \\
\bottomrule
\end{tabular}
\end{table}


\begin{table}[t]
\caption{Primary replacement summary. This table separates the main deployment comparison from exploratory rescue variants. Offline AUC is the covered-species OOF endpoint; private AUC is the hidden-test deployment endpoint.}
\label{tab:primary_summary}
\begin{tabular}{lccc}
\toprule
System & Offline AUC & Private AUC & Role \\
\midrule
ProtoSSMv2          & .803 & .942 & deployed reference \\
ProtoSSMv2 no-SSM   & --   & .941 & temporal-module ablation \\
Drop (Perch)        & .739 & .939 & no-prototype baseline \\
D1 pure             & .606 & .938 & high-purity replacement \\
D2 input-gated      & .830 & .937 & high-offline-AUC replacement \\
\bottomrule
\end{tabular}
\end{table}

\begin{figure}[t]
\centering
\includegraphics[width=\columnwidth]{fig_oof_vs_lb.pdf}
\caption{Offline OOF AUC (x-axis) vs.\ private leaderboard (y-axis) for variants with an offline OOF endpoint. Higher offline accuracy does not produce a higher deployment score: D2 (highest offline AUC, 0.830) is among the lowest on the leaderboard, while the deployed ProtoSSMv2 (middling offline AUC) is highest. The no-SSM temporal ablation has no separately computed offline endpoint and is therefore omitted from this scatter. The component ranking reverses direction (the absolute gaps are small; see Section~\ref{sec:limitations}).}
\label{fig:inversion}
\end{figure}

Table~\ref{tab:leaderboard} reports the hidden-test replacement results, and Table~\ref{tab:primary_summary} isolates the primary comparison. Each row corresponds to one submitted system generated by replacing or ablating the prototype-head artifact while leaving the rest of the Proto+SED core unchanged. Hidden-test labels are unavailable, so we cannot compute per-species confidence intervals or paired statistical tests for these scores. The leaderboard results are therefore descriptive system-level outcomes.

The primary comparison is between the deployed ProtoSSMv2 member, the two primary replacement heads (D1 pure and D2 input-gated), and the raw-Perch Drop baseline. The deployed ProtoSSMv2 member has the highest private score (0.942). The raw-Perch Drop baseline scores 0.939. The two primary replacement heads score below the Drop baseline: D1 scores 0.938 and D2 scores 0.937.

The offline ranking does not transfer to deployment. D2 has the highest offline AUC among the replacement heads (0.830), but it is the lowest primary replacement on the hidden test. D1 has much lower offline AUC (0.606) but slightly higher private score than D2 (0.938 vs. 0.937). These gaps are small and rest on a single hidden test (Section~\ref{sec:limitations}), but the direction is consistent across public and private splits.

The exploratory rescue variants recover only part of the gap. Coverage-gating and re-weighting cluster near the raw-Perch floor (0.938--0.940), and none reaches the deployed ProtoSSMv2 score. The temporal-module ablation provides a separate diagnostic: removing the deployed head's temporal SSM lowers private AUC from 0.94218 to 0.94089, but the ablated deployed head still ranks above the raw-Perch Drop baseline and both primary replacement heads. Section~\ref{sec:mechanism} analyzes these patterns.

\begin{figure}[t]
\centering
\includegraphics[width=\columnwidth]{fig_lb_bars.pdf}
\caption{Private leaderboard scores for all ten system variants. The deployed ProtoSSMv2 outperforms every replacement, ablation, and rescue variant. The no-SSM ablation is second-best, indicating that temporal aggregation contributes to the deployed head but does not explain its full advantage. D2-FOCAL (broad-trained on focal+soundscape data) scores worst. The y-axis is zoomed to show ordering; absolute score gaps are small.}
\label{fig:lbbars}
\end{figure}


%% ============================================================
%%  8. MECHANISM
%% ============================================================
\section{Mechanism Decomposition}
\label{sec:mechanism}

The offline-to-deployment inversion requires explanation.
We decompose it through an absent-species diagnostic, coverage-gating controls, a temporal-module ablation, and a broad-training sensitivity experiment.

\subsection{Diagnostic: the absent-species problem}

\begin{table}[t]
\caption{D2 gate and score diagnostics, split by species coverage. D2 dampens Perch's discriminative signal even on the 163 species it never trained on.}
\label{tab:mechanism}
\begin{tabular}{lrrr}
\toprule
Species group & Mean $\gamma$ & $\rho$(D2, Perch) & Var.\ ratio \\
\midrule
Covered (71)      & 0.779 & 0.319 & 0.485 \\
Absent (163)      & 0.767 & 0.276 & 0.347 \\
\bottomrule
\end{tabular}
\end{table}

Our training bank covers 71 of 234 species.
Table~\ref{tab:mechanism} shows that D2's input-dependent gate stays at $\gamma \approx 0.77$ even on the 163 absent species, reducing Perch's discriminative scores with a flat, untrained prototype signal. Correlations and variance ratios are computed over species within each coverage group $G$; for D2 scores $s^{\mathrm{D2}}$ and Perch scores $s^{\mathrm{Perch}}$, we report
\[
  \operatorname{VarRatio}_G =
  \frac{\operatorname{Var}_{c\in G}(s_c^{\mathrm{D2}})}
       {\operatorname{Var}_{c\in G}(s_c^{\mathrm{Perch}})}.
\]
Only approximately 35\% of Perch's score variance survives for absent species (vs.\ 49\% for covered species). Because macro-AUC averages over evaluable species in the hidden test, degradation across many target classes can affect the deployment score even when those classes are absent from the offline bank.

\subsection{Control 1: Coverage-gating}

We applied coverage-gating, using the prototype head for covered species and raw Perch logits for absent ones:
\[
S_c =
\begin{cases}
S_c^{\mathrm{proto}}, & c \in \mathcal{C}_{\mathrm{covered}},\\
S_c^{\mathrm{Perch}}, & c \notin \mathcal{C}_{\mathrm{covered}}.
\end{cases}
\]
This improved both heads on both splits:
D1 rose from 0.938 to 0.940 ($+0.002$), D2 from 0.937 to 0.938 ($+0.001$).
This pattern is consistent with absent-species degradation contributing to the deployment gap.
However, both heads recovered only to the Perch floor ($\approx$0.939), the score achieved by dropping the prototype member entirely.
Coverage-gating is necessary but not sufficient to match the deployed head.

\subsection{Control 2: Temporal aggregation explains part, not all, of the deployed-head advantage}

Because the deployed prototype member uses a temporal SSM while the primary replacement heads are window-level models, we also ablated the deployed head's temporal module. Removing the bidirectional SSM and associated temporal aggregation reduced private macro-AUC from 0.94218 to 0.94089 (Table~\ref{tab:leaderboard}). Thus, temporal aggregation contributes measurably to the deployed head's score.

However, the no-SSM ablation remains above the raw-Perch Drop baseline (0.941 vs.\ 0.939) and above both primary replacement heads (0.938 and 0.937). Temporal modeling therefore explains part of the deployed head's edge, but not all of it. The remaining advantage must come from non-temporal architectural or training differences, including how the deployed head preserves foundation-model signal across the full label space.

\subsection{Control 3: The deployed head preserves signal on absent species}

Coverage-gating the deployed ProtoSSMv2 itself \emph{hurt} it: $0.942 \rightarrow 0.940$.
This is unexpected: restricting the prototype to covered species should help if absent species are handled poorly.

Decomposing ProtoSSMv2's $+0.003$ edge over raw Perch:
only $+0.0003$ comes from covered species;
$+0.0026$ comes from absent ones.
Approximately 90\% of the deployed head's advantage over raw Perch is concentrated on the species it never saw training positives for: there it preserves more of Perch's discriminative signal, whereas our linear heads dampen it, the opposite pattern.
Together with the no-SSM ablation, this result suggests a multi-causal gap. Temporal aggregation accounts for a measurable part of the deployed head's advantage, but absent-species signal preservation remains after the temporal module is removed. The comparison to our lower-capacity heads remains confounded by architecture, capacity, and training procedure, so we do not attribute the reversal to temporal aggregation alone (Section~\ref{sec:limitations}).

\subsection{Capstone: Broad training backfires}

We next tested whether broader species coverage during training improves deployment performance.
We built D2-FOCAL: the gated head trained on the combined bank, the 708 soundscape windows plus focal recordings spanning 206 species, broadening coverage well beyond the 71 soundscape-labelled species.
We first verified that our local Perch extraction reproduces the deployed one exactly (embedding correlation 1.000000 on the 708-window bank), then embedded focal recordings across 206 species (5,090 windows, top-2 per species by logit confidence) and trained on the combined bank.

D2-FOCAL scored worst of all variants: private 0.935.
The focal recordings are off-domain relative to the soundscape test set, and 5,090 focal windows swamped the 708 on-domain ones.
Prototypes became focal-centric, matching the deployment distribution less well.
Broadening training data across the focal$\rightarrow$soundscape domain boundary, with this particular focal-heavy recipe, reduced performance. We tested only one recipe; balanced focal/soundscape sampling, domain weighting, or pseudo-labeling on the target distribution might recover the gap and remain untested.

\subsection{Summary}

The deployed head's advantage is consistent with multiple small factors. Coverage-gating shows that absent-species degradation contributes to the deployment gap, while the no-SSM ablation shows that temporal aggregation also contributes. Neither factor alone explains the full difference: coverage-gating recovers only to the Perch floor, and the no-SSM deployed head still beats both primary replacements. The tested broad off-domain training recipe reduces performance rather than closing the gap.


%% ============================================================
%%  9. DISCUSSION
%% ============================================================
%% ============================================================
%%  SUPPORTING STUDY (CONTROLLED SETTING)
%% ============================================================
\section{Supporting Study: Prototype Interpretability in a Controlled Setting}
\label{sec:validity}

The reversal above concerns whether an interpretable component \emph{transfers} to deployment, not whether prototypes can be valid at all. We checked the latter in a clean, single-backbone setting: a faithful AudioProtoPNet head~\cite{heinrich2025audioprotopnet} versus an identical sigmoid baseline on a shared ConvNeXt-Base backbone, 5-fold. The prototype head costs near-zero macro accuracy (macro $\Delta = -0.007$, 95\% CI spanning zero), with only a small, robust avian cost ($-0.011$). Its prototypes are class-pure for birds (0.951 versus a 0.006 null) and \emph{call-focused} rather than background-matching: their nearest patches avoid the low-energy tercile at every tested cutoff (a shortcut-rate test). So in a controlled setting prototype interpretability is close to free and behaves sensibly, which is precisely why the deployment reversal is notable: component validity in isolation did not predict deployment value. A compact version of this matched-control study is provided in Appendix~\ref{app:partone}.



\section{Discussion}
\label{sec:discussion}

\paragraph{Direct finding.}
The main observed result is an offline-to-deployment ranking failure under controlled replacement.
Inside the fixed ProtoSSMv2+SED core, the deployed prototype member remains best on the hidden private split, while both primary author-built replacement heads fall below the raw-Perch drop baseline.
This result is descriptive rather than inferential, because hidden-test labels are unavailable for paired testing, but it is the central empirical observation of the study.

\paragraph{Interpretation: partial-label evaluation exposes unmeasured failure modes.}
The proposed explanation is that the offline bank evaluates only the 71 species with positives in the labelled soundscape bank, while the deployment system must rank the broader 234-species target space.
The mechanism diagnostics support this interpretation in two ways.
First, D2 dampens Perch score variation more strongly on species absent from the offline bank than on covered species.
Second, coverage-gating the author-built heads improves hidden-test performance, indicating that some deployment loss comes from using prototype scores for species with no offline positives.
These findings do not prove that label coverage is the only mechanism; rather, they show that partial-label evaluation can hide damage on classes it never evaluates.

\paragraph{Temporal aggregation and architecture are partial contributors.}
The no-SSM ablation sharpens the mechanism claim.
Removing the deployed head's temporal module reduces private macro-AUC from 0.94218 to 0.94089, so temporal aggregation contributes to the deployed head's advantage.
However, the ablated deployed head still outranks both primary replacement heads and the raw-Perch drop baseline.
The deployment gap is therefore multi-causal: coverage handling, temporal aggregation, and remaining architectural differences each contribute to the final ranking.
The present evidence supports partial-label blindness as a failure mode, but it does not isolate every architectural factor in ProtoSSMv2.

\paragraph{Interpretable components must be evaluated as deployed components.}
The matched-control study (Appendix~\ref{app:partone}) shows that prototype explanations can be class-pure and call-focused in a cleaner single-backbone setting.
That result narrows the interpretation of the main failure: the problem is not that prototype explanations are inherently invalid, but that a locally interpretable component can still reduce deployment performance when it is selected using an incomplete offline label space.
For systems like BirdCLEF+ 2026, interpretability and deployment value are therefore separate endpoints.
A component can be more interpretable, or better on the covered offline bank, without improving the full system.

\paragraph{Taxon gradients are suggestive, not causal evidence.}
AudioProtoPNet's bird-centric pre-training produces a bird-clean purity gradient, while ProtoSSMv2's soundscape-trained Perch-based prototypes are cleaner for insects.
This contrast is consistent with the idea that prototype validity depends on training domain and acoustic distinctiveness.
However, the comparison is confounded by backbone, embedding space, training data, and prototype granularity.
We therefore treat the opposite taxon gradients as hypothesis-generating rather than as evidence for a causal taxon- or domain-level mechanism.

\paragraph{Practical recommendations.}
Three recommendations follow directly from the observed replacement results.
First, report label coverage for any offline component benchmark; otherwise, offline gains may be driven by a subset of classes that does not represent deployment.
Second, evaluate interpretable components inside the full system at least once, because component-level metrics do not capture interactions with rank fusion, foundation-model logits, or post-processing.
Third, when a component is trained on a subset of the deployment label space, include absent-class diagnostics or coverage-gated controls to check whether the component damages classes the offline metric cannot score.
These recommendations are scoped to BirdCLEF-style multi-label bioacoustic systems and should be tested in additional pipelines.

\paragraph{Negative results.}
Several negative results constrain the interpretation of the study.
The multi-label-aware separation loss performs similarly to standard separation.
ProtoSSMv2's per-class gate remains nearly uniform rather than learning an obvious coverage-aware allocation.
Increasing prototype count improves neither pure-head accuracy nor deployment performance.
Broad focal-heavy training worsens deployment in the tested recipe.
Finally, the initial expectation that the high-offline-AUC D2 head would transfer best to the leaderboard does not hold.
Together, these negative results suggest that the failure is not fixed by a single local design change.


%% ============================================================
%%  10. LIMITATIONS
%% ============================================================
\section{Limitations}
\label{sec:limitations}

\subsection{Statistical conclusion validity}

The offline evaluation uses 708 labelled soundscape windows from 59 files and covers only 71 of 234 target species.
This limits precision: fold-level uncertainty is wide, and the offline D2--ProtoSSMv2 difference should be interpreted descriptively rather than as a statistically secure ordering.
The hidden-test replacement results are also descriptive.
Kaggle returns one aggregate public score and one aggregate private score per submission, and hidden labels are unavailable, so we cannot compute paired tests, per-species confidence intervals, or submission-level p-values.
The public and private splits are not independent replications; they are two partitions of the same hidden evaluation source.
For these reasons, we present the leaderboard reversal as a documented case study in this pipeline rather than as a statistically established universal effect.

\subsection{Internal validity}

The controlled replacement design holds the ProtoSSMv2+SED core fixed and changes only the prototype member artifact, which reduces system-level confounding.
However, the replacement heads and the deployed ProtoSSMv2 member still differ in architecture, capacity, training procedure, and temporal structure.
The no-SSM ablation partially isolates temporal aggregation and shows that it contributes about 0.0013 private macro-AUC, but it does not isolate all non-temporal architectural differences.
The ablated deployed head still outperforms the raw-Perch drop baseline and both primary replacement heads, so the remaining advantage is confounded with capacity, fusion design, and other implementation choices.
A fully matched causal test would require training heads with matched capacity and identical temporal aggregation under the same data regime.

\subsection{Construct validity}

Presence-purity is a weak proxy for interpretability in multi-label soundscapes.
It verifies only that the nearest window contains the assigned class; it does not prove that the prototype matched the target call rather than a co-occurring species, background texture, recording condition, or correlated acoustic event.
The matched-control AudioProtoPNet study uses a stronger shortcut-rate check based on call energy, but that analysis is performed in a cleaner supporting setting rather than in the full deployed ProtoSSMv2 system.
Future interpretability validation should pair presence-purity with nearest-neighbor spectrogram panels, audio inspection, energy-overlap checks, and ideally expert review.

\subsection{External validity}

The results come from one BirdCLEF+ 2026 deployment pipeline and one competition hidden test.
The central lesson is therefore scoped: partial-label offline evaluation can mis-rank an interpretable component in this BirdCLEF-style system.
We do not claim that prototype heads generally fail, that temporal aggregation is always necessary, or that the same ranking reversal will occur in other bioacoustic datasets.
The broad-training result is also limited: D2-FOCAL used one focal-heavy recipe in which 5,090 focal windows outweighed 708 on-domain soundscape windows.
Balanced focal/soundscape sampling, domain weighting, pseudo-labeling on target soundscapes, and species-held-out validation remain untested.

\subsection{Supporting-study limitations}

The matched-control AudioProtoPNet study uses 6 training epochs, whereas the original AudioProtoPNet training used 20.
Longer training could reveal accuracy differences not captured here, especially for prototype regularization losses that may require more time to shape the latent space.
This limitation affects only the supporting AudioProtoPNet control and not the main ProtoSSMv2+SED replacement experiments.
The comparison between AudioProtoPNet and the Perch-based heads is also confounded by backbone, embedding space, training data, and prototype granularity.
The opposite taxon gradients are therefore suggestive, not causal evidence.

\subsection{Reproducibility and artifact limits}

The deployment testbed is adopted unchanged from public Kaggle notebooks~\cite{protossmv2}; only the replacement prototype heads and analyses are author-designed.
The repository linked below is intended to document the replacement-head code, evaluation notebooks, and deployment-substitution workflow.
However, hidden-test labels are not public, so leaderboard scores cannot be recomputed outside the Kaggle evaluation system.
Any future public release should therefore separate fully reproducible offline analyses from leaderboard rows that can only be verified through the original submission environment.


%% ============================================================
%%  FUTURE WORK
%% ============================================================
\section{Future Work}
\label{sec:future}

The next step is not simply to add more replacement heads, but to isolate the mechanisms identified above in a prioritized sequence.
First, a matched-capacity temporal ablation should compare four conditions: deployed ProtoSSMv2, ProtoSSMv2 without temporal aggregation, an author-built head with matched capacity but no temporal aggregation, and the same author-built head with the deployed temporal aggregator.
This would separate temporal structure from capacity and fusion design.

Second, broader training should be revisited with controlled domain balance rather than the focal-heavy D2-FOCAL recipe.
The most direct follow-up is a small grid over focal/soundscape sampling ratios, with species-held-out validation and a target-soundscape pseudo-labeling condition.
This would test whether broader species coverage helps once focal recordings no longer swamp on-domain soundscapes.

Third, interpretability validation should move beyond presence-purity.
For each prototype family, nearest-neighbor spectrogram and audio panels should be paired with energy-overlap or call-localization checks, and a small expert-reviewed subset would help distinguish target-call matches from co-occurrence or background matches.

Fourth, the partial-label evaluation failure should be replicated in other BirdCLEF-style systems.
The key replication test is whether a component selected on a partially covered offline bank also damages absent or weakly covered classes after system-level deployment.
If the same pattern appears across backbones, teams, or years, the result would support a general evaluation warning rather than a single-pipeline case study.

%% ============================================================
%%  11. CONCLUSION
%% ============================================================
\section{Conclusion}
\label{sec:conclusion}

We used a fixed BirdCLEF+ 2026 ProtoSSMv2+SED core with private macro-AUC 0.942 as a deployment testbed and replaced its prototype member with author-built prototype heads.
A matched-control study (Appendix~\ref{app:partone}) confirms that faithful AudioProtoPNet provides interpretability at near-zero accuracy cost, with prototypes class-pure and call-focused for birds but degrading for more distant taxa.
Our author-built heads map an offline interpretability--accuracy frontier and, most importantly, show that this offline ranking fails to transfer to the real leaderboard, reversing direction consistently across the public and private splits (the absolute gaps are small; see Section~\ref{sec:limitations}).
A mechanism decomposition indicates that the reversal is associated with preservation of foundation-model signal on offline-absent species: the deployed head preserves more of this signal, while our linear heads dampen it. A no-SSM ablation shows that temporal aggregation contributes to the deployed head's advantage but does not fully explain it. Our one attempt to close the gap through broader, focal-heavy training reduced performance, consistent with a focal-to-soundscape domain mismatch.

The lesson is concrete and scoped: in this BirdCLEF-style system, a partial-label offline bank mis-ranked an interpretable component because deployment scored a broader label space, and the component's interaction with the rest of the pipeline amplified the effect.
Whether a prototype head is ``interpretable'' and whether it helps the deployed system are different questions, with opposite answers here.
We suggest that component evaluations for such systems report label coverage, behavior on absent classes, and at least one system-level replacement test, rather than relying on a partial offline ranking alone. This recommendation is based on a single pipeline case study and should be tested in additional BirdCLEF-style systems.




%% ============================================================
%%  CODE AVAILABILITY
%% ============================================================
\section*{Code Availability}

Supporting code and experiment artifacts for this working note are available at \url{https://github.com/Sonoform-Labs/BirdClef}. The repository documents the replacement-head training code, evaluation notebooks, figure-generation scripts, leaderboard-row metadata, and deployment-substitution workflow used for the offline tables and hidden-test replacement rows. Hidden-test labels are not public, so the leaderboard scores can be verified only through the original Kaggle evaluation environment.

%% ============================================================
%%  ACKNOWLEDGMENTS
%% ============================================================
\begin{acknowledgments}
We thank the BirdCLEF+ 2026 organizers (Stefan Kahl, Holger Klinck, and colleagues), the Kaggle platform, and the LifeCLEF lab for hosting the competition and encouraging working note submissions.
We acknowledge the public Kaggle notebook community, particularly hideyukizushi, whose ProtoSSMv2 notebook forms the backbone of the deployment testbed we use, for making high-quality baselines accessible to all participants.
We also thank the developers of the open-source tools used in this work: PyTorch~\cite{paszke2019pytorch}, Perch~\cite{ghani2023global}, BirdSet~\cite{rauch2024birdset}, and AudioProtoPNet~\cite{heinrich2025audioprotopnet}.
\end{acknowledgments}



%% ============================================================
%%  APPENDIX A: TRAINING DETAILS
%% ============================================================
\bibliography{references}

\appendix

%% ============================================================
%%  APPENDIX A: AUDIOPROTOPNET MATCHED-CONTROL STUDY
%% ============================================================
\section{AudioProtoPNet Matched-Control Study}
\label{app:partone}

We include a small matched-control study to check whether the main deployment failure should be interpreted as a general failure of prototype explanations. We compare an AudioProtoPNet-style prototype head, implemented following the official BirdSet repository,\footnote{\url{https://github.com/DBD-research-group/BirdSet}} against a sigmoid classification head on the same ConvNeXt-Base backbone pretrained on BirdSet-XCL~\cite{liu2022convnext,rauch2024birdset}. Both models use the same folds, input representation, optimizer, augmentations, and evaluation protocol; only the classification head differs.

The prototype head is accuracy-competitive with the sigmoid baseline. Across five folds, the macro-AUC difference is $-0.007 \pm 0.006$ (prototype minus sigmoid), with a 95\% confidence interval spanning zero. The only consistent cost appears for Aves ($-0.011 \pm 0.001$), while Amphibia and Insecta have noisy fold-level deltas that are not distinguishable from zero.

To test prototype validity, we project each prototype to its nearest patch in a clean single-species pool of focal recordings and single-species soundscape windows. We measure \textbf{purity}, the fraction of prototypes whose nearest patch belongs to their own class, and \textbf{shortcut rate}, the fraction whose nearest patch falls in the bottom-energy region. Avian prototypes are strongly class-pure (0.951 versus a 0.006 random-prototype null) and avoid low-energy shortcut patches. Amphibian prototypes are moderately pure (0.562) but less consistently call-focused. Insect validity is coverage-limited because only 3 of 28 insect species have sufficient clean clips for this measurement.

These results show that prototype explanations can be class-pure and call-focused in this cleaner single-backbone setting. They do not show that prototype explanations are generally valid in BirdCLEF+ deployment. Therefore, the main paper's leaderboard reversal should not be read as evidence that prototype interpretability is impossible in bioacoustics. Rather, it supports our narrower claim: prototype components that look valid offline may still fail when selected under partial-label evaluation and deployed inside a full BirdCLEF system.


\end{document}
